\section{Further research: specific questions and workable next steps}
\label{sec:research}

The results change the most useful next questions. Repeating searches
for the eleven settled identities would add little. The next stage
should seek additional functional relations, effective bounds, and
well-specified arithmetic statements. The following problems distinguish
a plausible method from the mathematical conclusion it would need
to establish.

\subsection{First priority: the remaining Gaussian bridge}

\paragraph{1. Prove the \(S_4\) relation with an exact colored certificate.}
After the two shuffle eliminations, Conjecture~\ref{conj:S4} is the
single concrete target
\[
7\Im\Li_{4,1}(i,-1)
=51g_{41}+24g_{32}+\frac{19\pi^5}{512}
-14\beta(4)\log2.
\]
This displayed equation is still conjectural. Its second argument
is \(-1\), so it is not one of the already reduced inverse-argument
values \(\Li_{a,b}(i,-i)\).
A promising approach is to work with iterated integrals on the alphabet
\(\{0,1,-1,i,-i\}\), include explicitly regularized shuffle and stuffle
relations, and track the additional transformations of the punctured
line. The desired output is a rational linear combination of justified
relations with all boundary constants recorded.
An exact certificate would be a substantial improvement over another
decimal match.

\paragraph{2. Audit the colored relation system before enlarging the
constant basket.}
The inherited mixed formulas have modest integer coefficients, yet the
current manuscript reports them as apparent non-reductions.
Reproduce that computation from its value-generation stage, keeping
argument ordering and the parity component explicit.
At each weight record the exact alphabet, admitted product rows,
regularization prescriptions, distribution identities, and the
coefficient-height limit used by any numerical search.
Then compare the resulting exact row space with the explicit formulas
in Section~\ref{sec:mixed-reconciliation}.
This should determine whether the discrepancy comes from the numerical
values, a missing basis element, or an incomplete relation system.
An unsuccessful finite search should never be promoted to a statement
of period independence.

\paragraph{3. Determine what survives at weight six and depth three.}
The classes of \(F_{4,1,1},F_{3,2,1},F_{3,1,2}\) are an exact basis
of the formal one-variable shuffle quotient.
At \(i\), additional colored functional equations can reduce their
span further. The next question is the rank after specifying those
additional relations and the allowed lower-depth algebra.
The all-depth one-two theorem does not answer this: these indices
contain three zero letters in their words.
A useful outcome would be either explicit reductions or a certified
presentation of a larger, precisely defined formal quotient.
Neither outcome by itself proves that the numerical values are
linearly independent.

\subsection{Algebraic ladders beyond the two local proofs}

\paragraph{4. Produce local proofs of the higher golden ladders.}
The manuscript's higher golden formulas should receive the same kind
of argument/substitution and constant-term certificate as the cubic
identity. Start with the lowest unresolved weight and its full
functional-equation closure.
A vanishing symbol or differential can narrow the problem, but the
remaining products, zeta terms, and endpoint constants still require
determination. In particular, an identity valid only at an isolated
algebraic number cannot simply be integrated in that number to obtain
a higher-weight identity. A valid deformation must preserve the
functional relations being used.
Success here means a branch-consistent analytic proof of one exact
recorded formula, not only a lower-height numerical relation.

\paragraph{5. Classify useful complementary-power specializations.}
Theorem~\ref{thm:complementary-powers} starts from any
\(\rho^a+\rho^b=1\), but different exponent pairs may describe the
same field or the same base after a power substitution.
For bounded degree and exponent height, classify these redundancies
and determine which Kummer substitutions close on a small set of
powers of \(\rho\).
The cubic proof suggests looking for a closure that introduces a few
auxiliary powers and then cancels them through short rational rows.
The output can be an exact finite database of argument identities
over number fields, together with the corresponding polylogarithm
certificates. This is a constructive classification problem even when
period independence is out of reach.

\paragraph{6. Develop a reusable algebraic functional-equation verifier.}
Generalize the current quotient-ring calculation to arbitrary
irreducible polynomials, with an isolating interval or another exact
embedding specification. The verifier should check denominator
nonvanishing, identify the real or complex embedding, and retain
sign information needed for ordinary logarithms.
For single-valued functions, argument reduction is often sufficient;
for ordinary polylogarithms, continuation paths and inversion
corrections must remain part of the certificate.
The natural first benchmark is to reproduce the cubic certificate
without any special-case algebra in the program.

\subsection{Quantitative analysis of log-gamma moments}

\paragraph{7. Find the signed remainder near the least term.}
Theorem~\ref{thm:effective-moment} gives an unsigned bound of scale
\((e\alpha/n)^n\), within a polynomial factor of the smallest term.
It does not determine the leading signed error or prove an optimal
constant. Study \(R_{n,K}\) for
\(K=(n+3/2)/\alpha+O(1)\), allowing the integer rounding and the
alternation in \(k\) to enter explicitly.
The exact inverse-function integral is a starting point; the current
positive majorant discards cancellations that may control a sharper
answer. Any proposed asymptotic should first be checked over both
parities of \(K\), then proved uniformly in the rounding parameter.

\paragraph{8. Extend the theorem to weighted and reflected moments.}
For \(\beta>-1\), consider
\[
M_n(\beta)=\int_0^1x^\beta\log^n\Gamma(x)\,dx.
\]
The same layer-cake argument gives the exact starting identity
\[
\frac{M_n(\beta)}{n!}
=\frac{1}{(\beta+1)\Gamma(n)}
\int_0^\infty y^{n-1}x(y)^{\beta+1}\,dy,\qquad n\ge1.
\]
This identity follows by integrating \(x^\beta\) over the sublevel
interval \(0<x<x(y)\); it is not a conjecture.
For integer \(\beta+1\), Lagrange--B\"urmann coefficients of powers
of \(h\) suggest a direct extension of the coefficient analysis.
For noninteger powers, fractional leading exponents and branch
normalizations must be handled explicitly.
Mixed reflected moments involving
\(\log\Gamma(x)\) and \(\log\Gamma(1-x)\) add a second endpoint
mechanism and should be treated separately.

\paragraph{9. Determine the next inverse-gamma singularities.}
The local-limit proof finds the dominant growth without assuming a
global inverse-branch picture. A further analytic study should determine
the critical values of the reciprocal-gamma map accessible from this
branch, their local types, and the continuation paths that reach them.
That information could organize subsequent exponentially small
corrections and explain the behavior of higher-order truncation.
Computing roots of the digamma function numerically is a useful
exploration, but is not a proof of completeness or of which critical
value controls a given branch.

\paragraph{10. Reduce the known Tornheim coordinates in a specified
comparison algebra.}
The existence of a named Tornheim representation for \(M_3\) is settled
by \cite{BBB}. A meaningful remaining problem is, for example, whether
the particular combination \(2T_{13}-T_{12}\) at \((1,1,1)\) admits
a reduction using an explicitly stated set of zeta derivatives,
logarithms, and other standard constants.
The set must be fixed before a numerical search, and failure to
find a relation must retain its precision and height qualifications.
Independently, improve the triangular tail estimate
\eqref{eq:tornheim-tail} by summing in \(q=m+n\) and applying
an Euler--Maclaurin expansion to the resulting harmonic-logarithmic
coefficients. This could certify individual Tornheim combinations
without using the already evaluated cubic integral.

\subsection{A longer-term exactness programme}

\paragraph{11. Separate functional reduction from arithmetic
independence.}
The present results supply upper bounds and explicit reductions.
Lower bounds for numerical period spaces require different input.
One direction is to formulate the same questions in a motivic or
other formal period algebra, where a structural dimension can
sometimes be established, and then state exactly what comparison
or period conjecture would transfer that dimension to numbers.
The formal shuffle quotient of Theorem~\ref{thm:witt} is a useful
combinatorial baseline, but it contains too few relations to serve
as that arithmetic model on its own.

\paragraph{12. Move the smallest certificates into a proof assistant.}
The first targets should be the finite word-shuffle normal forms,
the binomial matrix inverse, and the rational algebraic substitutions.
They need only combinatorics and exact arithmetic.
The next layer consists of analytic lemmas: convergence of the
boundary Fourier convolution, the stated branches of polylogarithm
inversion, and a rigorous positive-tail implementation for the cubic
moment. A machine-checked rational row is valuable only when the
functional equations it combines have been justified in the same
convention. This staged approach keeps the trusted analytic base
visible and makes each new certificate reusable.

The common objective is a research record in which a numerical
discovery can be followed by a short exact witness, and in which the
questions that remain are stated at the level actually supported
by the evidence. The two-product shuffle, the algebraic ladder rows,
and the inverse-gamma remainder estimate provide three concrete
models for that next stage.
